% Method overview. Hand-written: it names parts and roles, not numbers.
\begin{table}[!htbp]
\caption{\textbf{What a view stores, and what each part is for.} A view is one encoder with the network trained on its rankings, and ArrowFlow trains $K$ of them (Section~\ref{sec:encoding}). A target-aware view puts coordinates fitted on the class labels before random ones, and a calibrated view standardizes its projected coordinates (Section~\ref{sec:encoder}). The output layer serves training and makes none of ArrowFlow's predictions.}
\label{tab:overview}
\centering
\footnotesize
\begin{tabular}{@{}p{0.19\textwidth} p{0.37\textwidth} p{0.34\textwidth}@{}}
\toprule
Part & What is fitted or stored & Role \\
\midrule
Encoder & imputation and standardization statistics and the projection matrix; a target-aware view also fits linear discriminant coordinates on the training labels, and a calibrated view stores the means and scales of its projected coordinates & turns a real-valued example into a ranking \\
\addlinespace
Hidden ranking layers & one permutation per ranking filter, reordered by position votes & produce the ranking that prediction reads \\
\addlinespace
Output layer & one permutation per class & gates the update, returns the signal from which every hidden layer learns, and scores the checkpoint \\
\addlinespace
Nearest-neighbor readout & the last hidden ranking of every training row with its label, and the chosen neighbor count and weighting & predicts the class of one view \\
\addlinespace
Vote across views & the predicted class of each view & gives the prediction of the ensemble \\
\bottomrule
\end{tabular}
\end{table}
